% ===================================================================
%  TABELL Nr. 12 — D'Gare. --- D'Eisebunnen. --- D'Schëffer.
%  Traduction 2026 d'apres texto/io, controlee sur texto/fr, texto/de
%  et texto/nl. Consigne : voir texto/lb/10-tabell-01.tex.
%
%  LE CHEMIN DE FER LUXEMBOURGEOIS PORTE LES DEUX COMPAGNIES QUI
%  L'ONT FAIT, ET LE TABLEAU 12 LES DONNE A LIRE MOT PAR MOT.
%
%  Est francais tout ce qui a trait au BATIMENT ET AU PAPIER :
%  « Gare », « Quai », « Guichet », « Billjee », « Buffet »,
%  « Compartiment », « Fourgon », « Affiche ». Est germanique tout ce
%  qui a trait a LA MACHINE ET A LA VOIE : « Eisebunn », « Zuch »,
%  « Schinnen », « Lokomotiv », « Kessel », « Kolben », « Rieder »,
%  « Damp », « Heizer », « Weichesteller », « Schlofwagon ».
%
%  ET CE PARTAGE N'EST PAS UN HASARD DE VOCABULAIRE : c'est
%  l'histoire de la ligne. Le reseau a ete construit sous concession
%  francaise et exploite ensuite sous administration allemande ; les
%  guichets et les billets sont restes du premier, les machines et
%  les voies du second. LE TABLEAU 11 AVAIT MONTRE QU'ON LIT DANS LE
%  NOM D'UN BUREAU DE QUEL COTE DE LA FRONTIERE IL EST ENTRE ; LE 12
%  MONTRE QU'ON PEUT LE LIRE DANS UN SEUL BATIMENT, ET MEME DANS UN
%  SEUL TRAIN.
%
%  LA COLONNE LITUANIENNE AVAIT TROUVE ICI UNE AUTRE LIGNE : tout le
%  mouvement nomme, tout le materiel emprunte. Elle ne joue pas ici,
%  parce que le luxembourgeois n'a rien nomme du tout — il a pris les
%  deux moities toutes faites, chacune a son fournisseur. Une langue
%  qui a deux voisins peut se dispenser d'inventer.
% ===================================================================

%%K t12-apar-1 apar
\VUsaut{4.0mm}
\VUtitre{15.0pt}{60}{TABELL Nr. 12}
\VUsaut{2.4mm}
\VUpk{11.4pt}{40}{D'Gare. --- D'Eisebunnen. --- D'Schëffer.}
\VUsaut{3.4mm}
\textsc{Bemierkung.} --- \textit{D'Fahrpläng, déi an all Natioun benotzt ginn, si
verschidden, dofir huelen mir als Beispill d'Stonne vun der} \VUgras{franséischer
Tabell.}

%%K t12-01-1 p
1. --- Ech si grad duerch Däitschland gereest. Virun der Ofrees hunn ech e
\VUgras{Fahrplangsbuch} \textsuperscript{(15)} kaaft, fir d'Ofreesstonn vun den Zich
ze wëssen. Nodeems ech festgestallt hat, datt dee gënschtegsten Zuch owes um néng
Auer vu Paris fortfiert an um eng Auer dräizéng zu Stroossbuerg ukënnt, hunn ech meng
Rees virbereet, an den Dag drop hunn ech mech an d'Gare fuere gelooss.

%%K t12-02-1 p
2. --- Wéi ech an déi grouss Ofreeshal erakoum, hannert engem alen Duerf-\VUgras{Paschtouer}
\textsuperscript{(2)}, deen a senger Hand seng \VUgras{Reesetäsch}
\textsuperscript{(3)} dréit, waren scho vill \VUgras{Reesend} \textsuperscript{(1)}
ukomm.

%%K t12-02-2 p
Well ech eng \VUgras{Depesch (Telegramm)} \textsuperscript{(5)} wollt schécken, hunn
ech mer vun engem \VUgras{Kontrolleur} \textsuperscript{(6)} de
\VUgras{Telegrafenbüro} \textsuperscript{(4)} weise gelooss.

%%K t12-03-1 p
3. --- Ier ech an de \VUgras{Waardesall} \textsuperscript{(7)} gaange sinn, hunn ech
mer e \VUgras{Billjee} \textsuperscript{(9)} mat Retour geholl.

%%K t12-04-1 p
4. --- Ech hunn net laang um \VUgras{Guichet} \textsuperscript{(8)} gewaart, well do
waren nëmme wéineg Leit. E \VUgras{Zaldot} \textsuperscript{(10)}, deen op Congé
gefuer ass, war esou léif, mir seng Plaz ze loossen, sou datt ech genuch Zäit hat, e
puer Informatioune beim \VUgras{Garechef} \textsuperscript{(13)} ze froen, dee mam
Iwwerwaachungs-\VUgras{Kommissär} \textsuperscript{(12)} a mam \VUgras{Gendaarm}
\textsuperscript{(11)} am Déngscht geschwat huet.

%%K t12-tit-1 sub
\VUpk{10.2pt}{40}{D'Bagage-Aschreiwung.}
\VUsaut{3.0mm}

%%K t12-05-1 p
5. --- Nodeems ech eng Zeitung um \VUgras{Zeitungskiosk} \textsuperscript{(14)} kaaft
hunn, hunn ech meng \VUgras{Bagage} \textsuperscript{(17)} weie gelooss, déi e
\VUgras{Portier} \textsuperscript{(16)} grad mat enger \VUgras{Bagagekar}
\textsuperscript{(19)} bruecht hat; de \VUgras{Beamten} \textsuperscript{(22)}, dee
fir d'\VUgras{Wo} \textsuperscript{(21)} zoustänneg ass, huet mer matgedeelt, datt
meng Kofferen dräiverzeg Kilo weien; dat huet e \VUgras{Iwwergewiicht} vu fënnef Kilo
gemaach, fir dat ech e Supplement um \VUgras{Bagage-Guichet} \textsuperscript{(23)}
bezuele muss.

%%K t12-06-1 p
6. --- Well ech ze fréi war, hat ech Zäit, de Verkéier vun de Reesenden an der Gare
ze beobachten. Déi eng hunn hir kleng Bagage am \VUgras{Bagage-Depot}
\textsuperscript{(25)} ofginn oder erëmgeholl; déi aner hunn de \VUgras{Fahrplang}
\textsuperscript{(26)} gekuckt. D'Reesend, déi ukomm sinn, sinn op den
\VUgras{Ausgang} \textsuperscript{(32)} zougeeilt, mat hirer \VUgras{Valis}
\textsuperscript{(24)} an der Hand. E puer hunn an der \VUgras{Bagage-Ausgab}
\textsuperscript{(27)} gewaart an hire \VUgras{Billjee} \textsuperscript{(29)}
gewisen, fir hir Kofferen, \VUgras{Këschten} \textsuperscript{(28)} oder
\VUgras{Kierf} \textsuperscript{(30)} erëmzekréien.

%%K t12-07-1 p
7. --- D'\VUgras{Akzisebeamten} \textsuperscript{(33)} hunn d'Päck suergfälteg
gepréift, fir festzestellen, ob een eppes ze deklaréieren huet.

%%K t12-08-1 p
8. --- E puer Reesend sinn an de \VUgras{Buffet} \textsuperscript{(34)} gaangen, wou
e \VUgras{Garçon} \textsuperscript{(35)} sech Méi ginn huet, si ze bedéngen. Si musse
sech beeilen, well den \VUgras{Zuch} \textsuperscript{(36)}, deen en Express ass,
bleift net laang an der Gare.

%%K t12-09-1 p
9. --- Duerno sinn ech op de Ofreesquai gaangen an hunn en direkte \VUgras{Wagon}
\textsuperscript{(37)} gesicht, fir net gezwongen ze sinn, ënnerwee de Wagon ze
wiesselen. Ech sinn an e Korridorwagon geklommen, deen e \VUgras{Compartiment}
\textsuperscript{(38)} fir Fëmmert an ee fir „Damme eleng“ huet.

%%K t12-tit-2 sub
\VUpk{10.2pt}{40}{Ofrees an der Nuecht.}
\VUsaut{3.0mm}

%%K t12-10-1 p
10. --- Nodeems ech d'\VUgras{Trëttbrett} \textsuperscript{(40)} eropgeklomme sinn,
d'\VUgras{Dier} \textsuperscript{(39)} zougemaach, de \VUgras{Store}
\textsuperscript{(41)} opgerullt a meng kleng Bagage op d'\VUgras{Netz}
\textsuperscript{(43)} geluecht hunn, hunn ech mech op d'\VUgras{Bänk}
\textsuperscript{(42)} gesat. Wéi ech de \VUgras{Luuchtemann} \textsuperscript{(44)}
d'Lampen unzéinne gesinn hunn, hunn ech geduecht, datt ech eng Nuecht am Wagon
verbrénge misst, an ech hunn de Beamte geruff, dee fir d'Verlounung vun de
\VUgras{Këssen} \textsuperscript{(45)} zoustänneg ass, fir hien no engem ze froen.

%%K t12-11-1 p
11. --- De Kondukter ass komm, fir d'Billjeeën ze kontrolléieren. Ech hunn hie
gefrot, firwat mir nach net fortgefuer sinn, well et ass jo déi richteg Stonn. Hien
huet mer geäntwert, datt de \VUgras{Zuchchef} \textsuperscript{(46)} amgaangen ass,
Vëloen an de \VUgras{Fourgon} \textsuperscript{(47)} ze setzen. Ausserdeem waren nach
net all Fousswärmer (48) gewiesselt.

%%K t12-12-1 p
12. --- Mä mir wäerte geschwënn fortfueren, well de \VUgras{Heizer}
\textsuperscript{(50)} huet op sengem \VUgras{Tender} \textsuperscript{(49)} Kuelen
geholl, fir d'Feier vun der \VUgras{Lokomotiv} \textsuperscript{(54)} unzefeieren, an
de \VUgras{Mechaniker} \textsuperscript{(51)} huet nëmmen op d'Signal vum
\VUgras{Ënnerchef vun der Gare} \textsuperscript{(52)} gewaart, deen um
\VUgras{Quai} \textsuperscript{(53)} stoung.

%%K t12-13-1 p
13. --- De \VUgras{Kessel} \textsuperscript{(56)} vun der Lokomotiv huet dumpf
gedonnert; de \VUgras{Kamäin} \textsuperscript{(55)} huet Strëmm vu Rauch mat
\VUgras{Damp} \textsuperscript{(60)} gemëscht ausgespien. Eng schrëll \VUgras{Päif}
\textsuperscript{(59)} huet gepaff, an d'\VUgras{Kolben} \textsuperscript{(58)} hu
sech bewegt an d'\VUgras{Rieder} \textsuperscript{(57)} vun der schwéierer Maschinn
gedriwwen.

%%K t12-tit-3 sub
\VUsaut{3.0mm}
\VUpk{10.2pt}{40}{Ofrees!}
\VUsaut{3.0mm}

%%K t12-14-1 p
14. --- D'Geschwindegkeet vum Zuch, deen op de \VUgras{Schinnen}
\textsuperscript{(64)} leeft, ass ëmmer méi grouss ginn; d'\VUgras{Quaien}
\textsuperscript{(65)} sinn verschwonnen; mir sinn un den \VUgras{Toiletten}
\textsuperscript{(66)} laanschtgefuer, an och un engem Wagonzuch, deen um anere
\VUgras{Gleis} \textsuperscript{(63)} stoung: \VUgras{Schlofwagonen}
\textsuperscript{(67)}, e \VUgras{Restaurantswagon} \textsuperscript{(68)}, e
\VUgras{Postwagon} \textsuperscript{(69)}.

%%K t12-noto-1 noto
\VUnotes{143.2mm}{%
\textit{(*) Bemierkung.} --- \textit{Selbstverständlech gëtt d'Rees no der Grenz vu
virum Krich beschriwwen.}%
}

%%K t12-15-1 p
15. --- Duerno ass d'\VUgras{Guttsgare} \textsuperscript{(71)} viru eisen Ae
laanschtgezunn. Ënner den Hallen hunn d'Beamten d'\VUgras{Wueren}
\textsuperscript{(72)} gelueden, déi mam Ilgutzuch geschéckt ginn.
\VUgras{Rangéierer} \textsuperscript{(76)} hunn nieft dem
\VUgras{Waasserreservoir} \textsuperscript{(73)} \VUgras{Véiwagonen}
\textsuperscript{(75)} a \VUgras{Plattformwagonen} \textsuperscript{(79)}
rangéiert, fir e \VUgras{Guttszuch} \textsuperscript{(80)} zesummenzestellen.

%%K t12-16-1 p
16. --- Mir sinn iwwert déi lescht \VUgras{Weiche} \textsuperscript{(78)} gefuer,
nieft deenen de Weichesteller stoung; mir sinn um \VUgras{Signaldisk}
\textsuperscript{(86)} laanschtgefuer, an d'Gare ass wäit ewech verschwonnen.

%%K t12-17-1 p
17. --- Geschwënn si mir iwwert eng \VUgras{Eisebréck} \textsuperscript{(74)}
gefuer, déi iwwert de \VUgras{Floss} \textsuperscript{(81)} geet. Nodeems mir déi
Bréck iwwerfuer haten, si mir laanscht de Floss gefuer, op deem mächteg
\VUgras{Schlepper} \textsuperscript{(82)} schwéier \VUgras{Lastbooter}
\textsuperscript{(83)} gezunn hunn. Mir hunn nach e \VUgras{Passagéiersschëff}
\textsuperscript{(84)} gesinn, wéi et un engem \VUgras{Ponton} \textsuperscript{(85)}
ugeluecht huet.

%%K t12-18-1 p
18. --- Nodeems mir ganz séier e puer Garen iwwerfuer sinn, si mir duerch e puer
\VUgras{Tunnelen} \textsuperscript{(87)} gefuer an hunn hannert eis onzieleg
\VUgras{Haischer} \textsuperscript{(88)} vun de \VUgras{Barrièrewärter} gelooss. Vum
Rëselen vum Zuch gewéit, si mir déif agepennt.

%%K t12-tit-4 sub
\VUsaut{3.0mm}
\VUpk{10.2pt}{40}{D'Gare vun Deutsch-Avricourt.}
\VUsaut{3.0mm}

%%K t12-19-1 p
19. --- Ech gouf op eemol vun der Stëmm vun engem Beamte gewächt, dee gejaut huet:
„Deutsch-Avricourt! \textsuperscript{(*)}. All Leit erof!“ Do war d'Kontroll vum Zoll
an all déi léideg Grenzformalitéiten. Ech hu missen meng Posch-Auer no der
\VUgras{Gareauer} \textsuperscript{(89)} stellen, déi fofzeg-a-fënnef Minutte
virgaangen ass; kaum wakereg, hunn ech de \VUgras{grousse Weiser}
\textsuperscript{(90)} nëmme schwéier vum \VUgras{klengen} \textsuperscript{(91)}
ënnerscheet; dat \VUgras{Zifferblat} \textsuperscript{(92)} selwer war wéinst dem
Moiesniwwel ganz onkloer.

%%K t12-20-1 p
20. --- Während d'Zollbeamten hir Kontroll gemaach hunn, hunn ech gäänzend
d'\VUgras{Affichen} \textsuperscript{(93)} entziffert, déi d'Garemaueren zéieren. Ech
hu gefriddstéckt, fir d'Zäit ze vertreiwen, hunn d'Kaart vum däitsche
\VUgras{Netz} \textsuperscript{(94)} gekuckt, fir ze gesinn, wéi vill Wee mech nach
vu Stroossbuerg trennt, a sinn nees a mäi Wagon geklommen, gestäerkt vun der
Hoffnung, geschwënn d'Zil vu menger Rees ze erreechen.

\VUsaut{6.0mm}
\VUfilet{22mm}
